% Response to Reviewer #1, Technical Issue (1)
% ISO 8608 comparison: What quantifiable advantages does the fractal framework offer?
%
% This file contains:
% 1. Analysis of why the issue arose
% 2. Assessment of legitimacy
% 3. Point-by-point response text
% 4. Proposed manuscript corrections

\documentclass[11pt]{article}
\usepackage[margin=1in]{geometry}
\usepackage{amsmath,amssymb}
\usepackage{xcolor}
\usepackage{hyperref}
\usepackage{setspace}
\onehalfspacing

\definecolor{revisionblue}{RGB}{0,0,180}

\begin{document}

\section*{Response to Reviewer \#1, Technical Issue (1)}

\subsection*{Reviewer Comment (verbatim)}

\begin{quote}
The manuscript establishes a fractal-based framework. However, a straightforward approach would be to directly generate road profiles from ISO 8608 and then use them as input to excite a quarter-car model for calculating vibration energy and fatigue damage. The authors should clarify what quantifiable advantages the proposed fractal framework offers over the direct use of ISO 8608 parameters, in terms of prediction accuracy, scope of applicability, or practical operability.
\end{quote}

%% ============================================================
\subsection*{1. Why This Issue Arose}
%% ============================================================

The reviewer's concern is well-founded. The original manuscript discussed the connection to ISO~8608 only briefly in the Discussion section (one paragraph noting that $G_d(n_0)$ corresponds to $C_z$ and the waviness exponent $w$ corresponds to $\beta$). This treatment was insufficient because:

\begin{itemize}
    \item It did not explicitly state \emph{what problem} the ISO~8608 direct approach has that motivates the fractal framework.
    \item It did not provide quantifiable metrics comparing the two approaches (e.g., prediction accuracy, number of parameters, applicability domain).
    \item The Introduction framed the gap as ``no quantitative framework connects terrain geometry to fatigue,'' but ISO~8608 \emph{does} provide a spectral characterization of roads---so the reader naturally asks why that existing standard is insufficient.
    \item The paper's own two-parameter model $(C_z, \beta)$ is structurally identical to ISO~8608's $(G_d(n_0), w)$, making the distinction unclear without explicit differentiation.
\end{itemize}

In short, the manuscript assumed the reader would recognize the limitations of ISO~8608 for natural/off-road terrain, but never articulated those limitations explicitly.

%% ============================================================
\subsection*{2. Is the Comment Legitimate?}
%% ============================================================

\textbf{Yes, fully legitimate.} This is a core novelty question that any reviewer would ask. The comment identifies a genuine gap in the manuscript's argumentation. Specifically:

\begin{enumerate}
    \item \textbf{ISO~8608 does provide a two-parameter PSD model} ($G_d(n_0)$ and waviness $w$), so the paper's $(C_z, \beta)$ framework is not structurally novel in parameterization---the novelty must lie elsewhere.
    
    \item \textbf{The direct simulation pipeline exists:} One can generate ISO~8608 profiles, run them through a quarter-car model, and compute fatigue. This is standard practice in automotive engineering. The paper must explain what it adds beyond this.
    
    \item \textbf{Three genuine advantages exist but were not articulated:}
    \begin{itemize}
        \item[(a)] \emph{Geometric interpretability:} ISO~8608 treats $w$ as an empirical constant ($w \approx 2$); the fractal framework derives it from terrain geometry ($\beta = 7-2D$), enabling prediction from DEMs without road measurements.
        \item[(b)] \emph{Scope beyond engineered roads:} ISO~8608 is calibrated for paved/graded roads (Classes A--H). It does not cover natural terrain where no road classification exists. The fractal framework applies to any terrain with a DEM.
        \item[(c)] \emph{Predictive capability without in-situ measurement:} ISO~8608 requires profilometer data to classify a road. The fractal framework can predict spectral parameters from remote sensing (satellite/airborne DEMs) before any vehicle traverses the terrain.
    \end{itemize}
\end{enumerate}

The reviewer is correct that these advantages must be stated explicitly and quantified where possible.

%% ============================================================
\subsection*{3. Point-by-Point Response (for submission)}
%% ============================================================

\noindent\textbf{Response:}

We thank the reviewer for this important observation. The reviewer correctly identifies that the original manuscript did not sufficiently differentiate the proposed framework from the direct ISO~8608 simulation pipeline. We have now added explicit discussion of the quantifiable advantages in both the Introduction and Discussion. The key distinctions are:

\begin{enumerate}
    \item \textbf{Geometric interpretability and predictability from remote sensing.} ISO~8608 treats the waviness exponent $w$ as a fixed empirical constant (typically $w \approx 2$) or as a parameter that must be measured via profilometer. Our framework derives this exponent from terrain fractal geometry ($\beta = 7-2D$), enabling prediction of spectral slope directly from digital elevation models (DEMs) without requiring in-situ road profile measurements. This is critical for military route planning, planetary rover operations, and any scenario where profilometer access is unavailable.

    \item \textbf{Scope of applicability.} ISO~8608 is calibrated for engineered road surfaces and classifies roads into discrete amplitude classes (A--H) based on $G_d(n_0)$ at a reference wavenumber. It does not extend to natural terrain (mountains, deserts, off-road environments) where no road classification exists. Our framework applies to any terrain for which a DEM is available, as validated by the USGS 3DEP LiDAR analysis ($r = -0.62$, $p < 0.05$, $n = 13$ terrain regions).

    \item \textbf{Two-parameter prediction accuracy.} While ISO~8608 traditionally uses a single amplitude parameter $G_d(n_0)$ for classification (assuming fixed $w \approx 2$), our results demonstrate that the two-parameter model $(C_z, \beta)$ improves vibration energy prediction from $R^2 = 0.91$ (amplitude only) to $R^2 = 0.96$ (both parameters). This 5\% improvement corresponds to a 53\% reduction in unexplained variance, which is physically meaningful for fatigue prediction where errors compound through the power-law S--N relationship (factor of $\sim$2 in predicted fatigue life for $m = 5$).

    \item \textbf{Physical causal chain.} The direct ISO~8608 pipeline treats the road profile as an empirical input with no connection to terrain geomorphology. Our framework establishes the causal chain from terrain geometry ($D$) through spectral structure ($\beta = 7-2D$) to vibration energy and fatigue, enabling understanding of \emph{why} certain terrains produce more damage and how geological processes create terrain severity.
\end{enumerate}

We emphasize that the proposed framework does not replace ISO~8608 for engineered roads where profilometer data are available. Rather, it extends the spectral approach to natural terrain and provides a geometric foundation for the waviness parameter. These distinctions are now explicitly stated in the revised manuscript (Introduction, paragraph 4; Discussion, expanded ISO~8608 paragraph). Changes are marked in {\color{revisionblue}blue}.

%% ============================================================
\subsection*{4. Manuscript Corrections}
%% ============================================================

Two insertions are made in the revised manuscript, marked with {\color{revisionblue}blue text} and marginal annotations [R1.1]:

\subsubsection*{Correction 1: New paragraph in Introduction (after the IRI paragraph, before the fractal analysis paragraph)}

{\color{revisionblue}
\textbf{[R1.1]} A direct approach to terrain-induced fatigue prediction would generate road profiles from the ISO~8608 spectral classification \cite{ISO8608} and use them as quarter-car model inputs. However, this pipeline has three limitations that motivate the present work. First, ISO~8608 treats the waviness exponent $w$ as a fixed empirical constant (typically $w \approx 2$) that must be measured via profilometer; it provides no mechanism to predict spectral slope from terrain geometry alone. Second, ISO~8608 is calibrated for engineered road surfaces (Classes A--H) and does not extend to natural terrain where no road classification exists---precisely the domain where fatigue prediction is most needed (military operations, autonomous off-road navigation, planetary exploration). Third, ISO~8608 classifies roads primarily by a single amplitude parameter $G_d(n_0)$, whereas our results demonstrate that the two-parameter description $(C_z, \beta)$ improves predictive fit from $R^2 = 0.91$ to $R^2 = 0.96$, a substantial reduction in residual variance that propagates through the nonlinear S--N relationship and can materially affect fatigue-life estimates. The fractal framework addresses all three limitations by deriving spectral slope from terrain geometry ($\beta = 7-2D$), extending applicability to any terrain with a digital elevation model, and quantifying the independent contributions of amplitude and spectral structure to fatigue damage.
}

\subsubsection*{Correction 2: Expanded paragraph in Discussion (replacing the existing ISO~8608 paragraph)}

{\color{revisionblue}
\textbf{[R1.1]} \textbf{Connection to ISO~8608 road roughness classification.} The proposed $(C_z, \beta)$ spectral framework is structurally related to the ISO~8608 road roughness classification, which models road PSD as $G_d(n) = G_d(n_0)(n/n_0)^{-w}$. In this formulation, $G_d(n_0)$ corresponds to the spectral amplitude $C_z$, while the waviness exponent $w$ corresponds to the spectral slope $\beta$. The present work extends this established framework in three quantifiable ways. (i)~\emph{Geometric origin:} We derive the waviness exponent from self-affine surface theory ($\beta = 7-2D$), providing a geometric basis for interpreting spectral slope in terrain PSDs. This explains why different terrain classes exhibit characteristic spectral slopes and enables prediction of $\beta$ from DEM-derived fractal dimension without profilometer measurements. (ii)~\emph{Scope extension:} ISO~8608 is calibrated for engineered roads where profilometer data are available. The fractal framework generalizes to natural terrain (validated on 13 USGS 3DEP regions spanning plains to mountains), enabling fatigue prediction from remote sensing data prior to vehicle deployment---a capability absent from the ISO~8608 pipeline. (iii)~\emph{Two-parameter improvement:} Standard ISO~8608 practice classifies roads by amplitude alone (assuming $w \approx 2$). Our analysis demonstrates that explicitly fitting both parameters improves predictive fit from $R^2 = 0.91$ to $R^2 = 0.96$, a substantial reduction in residual variance that propagates through the nonlinear S--N relationship and can materially affect fatigue-life estimates. The framework thus does not replace ISO~8608 for paved roads but provides a geometric interpretation of spectral slope within an analogous PSD framework and extends its applicability to uncharacterized natural terrain.
}

\end{document}
